% SPDX-License-Identifier: CC-BY-SA-4.0
% Author: Matthieu Perrin

\begingroup

\begin{exercice}[Mini-défi -- Étendre le lemme d'Arden]\label{exo:lexing/arden/challenge}

  Soient $\Sigma$ un alphabet et $A, B \subseteq \Sigma^\star$ deux langages. 
  On considère l'équation en langages $X = A\cdot X \mid B$.
  D'après le lemme d'Arden, nous savons que l'équation $X = A\cdot X\mid B$ a une seule
  solution quand $\varepsilon\notin A$.
  On suppose $\varepsilon \in A$. Le but de l'exercice est de prouver le théorème suivant :
  
  $$
  \text{Un langage }X\text{ est solution de }X = A\cdot X \mid B 
  \iff \exists L \subseteq \Sigma^\star,\; X = A^\star \cdot (L \mid B).
  $$

  \begin{question}
  \item Montrer que $\Sigma^\star$ est solution de $X = A\cdot X \mid B$.
  \item Montrer que, pour tout $L \subseteq \Sigma^\star$, le langage $A^\star \cdot (L \mid B)$ est solution de $X = A\cdot X \mid B$.
  \item Soit $M$ une solution de $X = A\cdot X \mid B$. Montrer que $M \subseteq A^\star \cdot (M \mid B)$.
  \item Montrer, par récurrence sur $n$, que $\forall n\in\mathbb{N},\; A^n \cdot (M \mid B) \subseteq M$.
  \item En déduire $M = A^\star \cdot (M \mid B)$.
  \item Conclure que $X$ est solution de $X = A\cdot X \mid B$ si et seulement si 
    $\exists L \subseteq \Sigma^\star$ tel que $X = A^\star \cdot (L \mid B)$.
  \end{question}

\end{exercice}

\begin{correction}

  On note $F(X) \eqdef A\cdot X \mid B$.

  \begin{question}

  \item \textbf{$\Sigma^\star$ est solution.}
    On a $A\cdot \Sigma^\star = \Sigma^\star$ et $B \subseteq \Sigma^\star$, donc
    $F(\Sigma^\star) = \Sigma^\star \mid B = \Sigma^\star$.

  \item \textbf{$A^\star(L \mid B)$ est toujours solution.}
    Comme $\varepsilon \in A$, on a $A\cdot A^\star = A^\star$.
    Alors, pour tout $L$ :
    \[
    F\big(A^\star(L \mid B)\big) 
    = A\cdot A^\star(L \mid B) \mid B
    = A^\star(L \mid B) \mid B
    = A^\star(L \mid B),
    \]
    car $B \subseteq A^\star(L \mid B)$.

  \item \textbf{$M \subseteq A^\star(M \mid B)$.}
    Comme $M \subseteq M \mid B$ et $\varepsilon \in A^\star$, on a $M \subseteq A^\star(M \mid B)$.

  \item \textbf{Récurrence sur $n$.}
    Initialisation ($n=0$) : $A^0(M \mid B) = M \mid B \subseteq M$ car $M = A\cdot M \mid B$ implique $B \subseteq M$.
    Hérédité : si $A^n(M \mid B) \subseteq M$, alors
    \[
    A^{n+1}(M \mid B) = A \cdot \big(A^n(M \mid B)\big) \subseteq A \cdot M \subseteq M,
    \]
    puisque $M = A\cdot M \mid B$.

  \item \textbf{$M = A^\star(M \mid B)$.}
    Par 3, $M \subseteq A^\star(M \mid B)$.
    Par 4, pour tout $n$, $A^n(M \mid B) \subseteq M$, donc $\bigcup_{n\ge 0} A^n(M \mid B) \subseteq M$.
    Or $A^\star(M \mid B)=\bigcup_{n\ge 0} A^n(M \mid B)$. D’où l’égalité.

  \item \textbf{Caractérisation.}
    ($\Rightarrow$) Si $X$ est solution, prendre $L \eqdef X$ et utiliser le point 5 appliqué à $M=X$ :
    $X = A^\star(X \mid B) = A^\star(L \mid B)$.

    ($\Leftarrow$) Réciproquement, le point 2 montre que tout $X$ de la forme $A^\star(L \mid B)$ est solution.

  \end{question}

\end{correction}

\endgroup
\endinput
